Consistent boundary conditions are a key input to large model intercomparison projects.
Here we present the greenhouse gas concentration boundary conditions, i.e. forcings, to be used in the seventh phase of the Coupled Model Intercomparison Project (CMIP7).
We show that, despite the use of updated input data sources, changes between the previous phase's forcings (CMIP6) and CMIP7 forcings are small:
CO_2 has a maximum deviation of [TODO X], with the maximum deviation since the beginning of ground-based observations in [TODO Y] being limited to [TODO X];
CH_4 has a maximum deviation of [TODO X], [TODO X] since [TODO Y];
N_2O has a maximum deviation of [TODO X], [TODO X] since [TODO Y]
and CFC12 [TODO check nomenclature] has a maximum deviation of [TODO X], since [TODO Y].
The forcing dataset also shows that, between 2014 (the end of CMIP6's historical dataset)
and 2022 (the end of these datasets),
concentrations of CO_2 have increased from [TODO X] to [TODO Y],
of CH_4 from [TODO X] to [TODO Y],
of N_2O from [TODO X] to [TODO Y]
while concentrations of CFC12 have fallen from [TODO X] to [TODO Y]/
In approximate effective radiative forcing terms (ERF), all changes are less than [TODO X]
and the biggest change between the two datasets is [TODO Y] W / m^2 in [TODO Z year].
These bigger changes are driven by revisions in the interpolating spline used over the period that is only covered by ice core based observations.
In the period covered by ground-based measurements, changes are much smaller.
These forcings are being used for CMIP7 historical and other simulations at the moment.
We conclude with a section on how to produce annual updates of these forcings,
to support more rapid evaluation of recent changes in future.
